% ============================================================
% Submission LaTeX source — Asia Pacific Journal of Management (Springer)
% Document class for final submission: Springer svjour3.cls (smallextended) — Springer Nature LaTeX support; article fallback để xác minh
% Compile: xelatex 04_manuscript_latex.tex (x2). Figures submitted separately.
% ============================================================
% Options for packages loaded elsewhere
\PassOptionsToPackage{unicode}{hyperref}
\PassOptionsToPackage{hyphens}{url}
\PassOptionsToPackage{dvipsnames,svgnames,x11names}{xcolor}
%
\documentclass[
  12pt,a4paper]{article}
\usepackage{amsmath,amssymb}
\usepackage{setspace}
\usepackage{iftex}
\ifPDFTeX
  \usepackage[T1]{fontenc}
  \usepackage[utf8]{inputenc}
  \usepackage{textcomp} % provide euro and other symbols
\else % if luatex or xetex
  \usepackage{unicode-math} % this also loads fontspec
  \defaultfontfeatures{Scale=MatchLowercase}
  \defaultfontfeatures[\rmfamily]{Ligatures=TeX,Scale=1}
\fi
\usepackage{lmodern}
\ifPDFTeX\else
  % xetex/luatex font selection
  \setmainfont[]{TeX Gyre Termes}
  \setmathfont[]{TeX Gyre Termes Math}
\fi
% Use upquote if available, for straight quotes in verbatim environments
\IfFileExists{upquote.sty}{\usepackage{upquote}}{}
\IfFileExists{microtype.sty}{% use microtype if available
  \usepackage[]{microtype}
  \UseMicrotypeSet[protrusion]{basicmath} % disable protrusion for tt fonts
}{}
\makeatletter
\@ifundefined{KOMAClassName}{% if non-KOMA class
  \IfFileExists{parskip.sty}{%
    \usepackage{parskip}
  }{% else
    \setlength{\parindent}{0pt}
    \setlength{\parskip}{6pt plus 2pt minus 1pt}}
}{% if KOMA class
  \KOMAoptions{parskip=half}}
\makeatother
\usepackage{xcolor}
\usepackage[margin=2.5cm]{geometry}
\usepackage{longtable,booktabs,array}
\usepackage{calc} % for calculating minipage widths
% Correct order of tables after \paragraph or \subparagraph
\usepackage{etoolbox}
\makeatletter
\patchcmd\longtable{\par}{\if@noskipsec\mbox{}\fi\par}{}{}
\makeatother
% Allow footnotes in longtable head/foot
\IfFileExists{footnotehyper.sty}{\usepackage{footnotehyper}}{\usepackage{footnote}}
\makesavenoteenv{longtable}
\setlength{\emergencystretch}{3em} % prevent overfull lines
\providecommand{\tightlist}{%
  \setlength{\itemsep}{0pt}\setlength{\parskip}{0pt}}
\setcounter{secnumdepth}{-\maxdimen} % remove section numbering
\ifLuaTeX
\usepackage[bidi=basic]{babel}
\else
\usepackage[bidi=default]{babel}
\fi
\babelprovide[main,import]{english}
\ifPDFTeX
\else
\babelfont{rm}[]{TeX Gyre Termes}
\fi
% get rid of language-specific shorthands (see #6817):
\let\LanguageShortHands\languageshorthands
\def\languageshorthands#1{}

\setcounter{secnumdepth}{-1}  % giữ numbering thủ công trong manuscript
\usepackage{booktabs}
\usepackage{longtable}
\usepackage{newunicodechar}
\newunicodechar{₀}{\textsubscript{0}}\newunicodechar{₁}{\textsubscript{1}}
\newunicodechar{₂}{\textsubscript{2}}\newunicodechar{₃}{\textsubscript{3}}
\newunicodechar{₄}{\textsubscript{4}}\newunicodechar{₅}{\textsubscript{5}}
\newunicodechar{≈}{$\approx$}\newunicodechar{≤}{$\leq$}\newunicodechar{≥}{$\geq$}
\newunicodechar{×}{$\times$}\newunicodechar{−}{$-$}
\ifLuaTeX
  \usepackage{selnolig}  % disable illegal ligatures
\fi
\IfFileExists{bookmark.sty}{\usepackage{bookmark}}{\usepackage{hyperref}}
\IfFileExists{xurl.sty}{\usepackage{xurl}}{} % add URL line breaks if available
\urlstyle{same}
\hypersetup{
  pdftitle={Does Country Context Shape the Internationalization--Performance Link? A Three-Level Meta-Analytic Investigation of Digital Adoption, Institutional Regimes, and the Digital Paradox Lifecycle},
  pdflang={en},
  colorlinks=true,
  linkcolor={black},
  filecolor={Maroon},
  citecolor={black},
  urlcolor={blue},
  pdfcreator={LaTeX via pandoc}}

\title{Does Country Context Shape the Internationalization--Performance
Link? A Three-Level Meta-Analytic Investigation of Digital Adoption,
Institutional Regimes, and the Digital Paradox Lifecycle}
\author{}
\date{}

\begin{document}
\maketitle
\begin{abstract}
We conduct a three-level meta-analytic regression analysis (MARA)
examining whether country-level digital adoption (cDAI), institutional
context regime (ICRV), and Digital Paradox Lifecycle (DPL) phase
moderate the internationalization--performance (I to P) relationship
three theoretically motivated moderators that prior meta-analyses have
not examined. A systematic search following PRISMA 2020 protocols on Web
of Science and Scopus (1977--2026) identifies \emph{k} = 238 studies
with \emph{K} = 288 effect sizes. Three-level MARA (Cheung, 2014; Van
den Noortgate et al., 2013) decomposes heterogeneity into within-study
(\emph{Level 2}) and between-study (\emph{Level 3}) components using
\texttt{metafor} in R. The analysis plan is registered on OSF
(transparency registration); inter-coder reliability κ ≥ 0.70 on 20\%
double-coded subsample. The baseline pooled effect is \emph{r} = 0.074
(95\% CI {[}0.060, 0.088{]}, \emph{p} \textless{} .001) with \emph{I}² =
62.4\% (within-study 54.1\%; between-study 8.4\%), replicating and
extending the ICBEF 2025 baseline (\emph{r} = 0.07, \emph{k} = 113). The
three hypothesized moderators show limited support: ICRV regime
differences are statistically significant (\emph{Q}\_M = 17.35,
\emph{df} = 4, \emph{p} = .002) but driven by an anomalous
Frontier-group estimate (\emph{k} = 3 studies, \emph{r̄} = 0.35) rather
than a monotone institutional gradient; cDAI and DPL phase moderation
are not statistically significant (\emph{Q}\_M = 1.23, \emph{p} = .541
and \emph{Q}\_M = 0.62, \emph{p} = .734, respectively) in the current
\emph{k} = 238 sample. Substantial publication bias is detected:
trim-and-fill imputes \emph{k} = 57 missing studies, reducing the
adjusted pooled effect to \emph{r} = 0.035 (95\% CI {[}0.018, 0.051{]}),
positive but attenuated. This is the first three-level MARA of I to P
drawing on \emph{k} = 238 studies from 49 economies globally, and the
first to formally test ICRV, cDAI, and DPL phase as moderators. The ICRV
between-regime Q\_M test confirms H1 (Q\_M = 17.35, \emph{p} = .002),
while the directional gradient (E1a/E1b) and the cDAI and DPL hypotheses
remain unconfirmed, a finding that, combined with substantial
publication bias, reframes the heterogeneity puzzle: unexplained
variance in I to P may reflect publication-side selection more than
institutional or digital contingencies, calling for pre-registered
replication with larger between-regime samples.
\end{abstract}

\setstretch{2.0}
\textbf{Keywords:} internationalization--performance; meta-analysis;
three-level model; digital adoption; institutional context; ICRV;
Digital Paradox Lifecycle; global; Asia-Pacific

\hypertarget{introduction}{%
\subsection{1. Introduction}\label{introduction}}

Cross-border firm expansion has become one of the defining strategic
imperatives of the modern economy. The past three decades have witnessed
dramatic growth in outward foreign direct investment, export-oriented
production, and global value chain participation, transforming
internationalization from a strategy available primarily to large
multinationals into a competitive consideration for firms across size
classes and geographies (Dunning, 2000; Johanson \& Vahlne, 2009;
Kafouros et al., 2012). Against this backdrop, the question of whether
internationalization improves firm performance is not merely academic;
it shapes investment decisions, government export promotion policy, and
the strategic priorities of firms navigating an increasingly
interconnected global economy (Hitt et al., 2006; Lu \& Beamish, 2004).
For scholars, however, the record remains stubbornly inconclusive.

The relationship between a firm's degree of internationalization and its
performance (the ``I to P relationship'') is the most meta-analyzed
question in international business (IB). Over four decades and six major
meta-analyses (Bausch \& Krist, 2007; Kirca et al., 2012; Marano et al.,
2016; Schwens et al., 2018; Wu et al., 2022; Arte \& Larimo, 2022), no
consensus has emerged: pooled effects are consistently small and
positive, yet \emph{I}² regularly exceeds 80\%, signaling that context,
not a universal mechanism, drives outcomes.

The present study's starting point is the ICBEF 2025 baseline analysis
(Authors, 2024): \emph{k} = 113 studies, pooled \emph{r} = 0.07
(\emph{p} \textless{} .001), \emph{I}² = 87.92\%. While confirming the
positive average effect, this baseline identified that conventional
moderators, country of origin, industry, performance measure type, leave
approximately 70\% of variance unexplained. Three theoretically grounded
moderators, absent from all prior meta-analyses, motivate the present
extension:

\textbf{Gap 1, cDAI.} Country-level digital adoption (cDAI) has been
proposed as a contextual amplifier of firm-level competitive advantages
(Stallkamp \& Schotter, 2021; Verhoef et al., 2021), yet no
meta-analysis has tested whether the national digital infrastructure
environment moderates the I to P link.

\textbf{Gap 2, ICRV 5-regime.} Marano et al.~(2016) established that
home-country institutions moderate I to P, but applied a coarse
six-group global taxonomy. The global study corpus spans the full
institutional spectrum, from Singapore (World Governance Indicators Rule
of Law score +1.84) to frontier economies such as Pakistan (WGI −0.55)
and Iran (WGI −0.74; World Bank, 2023). An ICRV 5-regime classification
capable of resolving this heterogeneity has not been tested
meta-analytically on a globally representative I to P corpus.

\textbf{Gap 3, DPL phase.} Brynjolfsson et al.~(2021) identified 2009 as
a productivity inflection point in the digital era (the ``dynamo
analogy'' for AI: David, 1990). Studies examining data from before,
spanning, or after this threshold should yield systematically different
I to P effect sizes if digital platforms reshape internationalization
economics. This temporal moderator has never been systematically coded
in I to P meta-analyses.

This paper addresses all three gaps through a fresh systematic search
expanding the ICBEF 2025 baseline from \emph{k} = 113 to \emph{k} = 238,
combined with three-level MARA that decomposes heterogeneity beyond what
random-effects models allow.

\textbf{Contributions.} We make three methodological and three
theoretical contributions: \emph{(Methodological)}: (1) First
three-level MARA for the I to P literature; (2) first
PRISMA-2020-compliant systematic search with OSF registration for this
topic; (3) first application of between-study \emph{vs.} within-study
heterogeneity decomposition to I to P. \emph{(Theoretical)}: (4) First
meta-analytic test of an ICRV 5-regime framework applied to a globally
representative I to P corpus (\emph{k} = 238, 49 economies); (5) first
formal test of cDAI as a national-level digital-infrastructure moderator
of I to P; (6) first test of DPL phase as a temporal moderator using
three-level MARA. The non-confirmation of E1a/E1b, H2, and H3, and the
anomalous ICRV Frontier pattern, are themselves informative findings
that bound the conditions under which these moderators could operate.

\textbf{Key findings preview.} The baseline pooled effect (\emph{r} =
0.074, \emph{k} = 238, \emph{K} = 288) replicates and extends the ICBEF
2025 finding (\emph{r} = 0.07, \emph{k} = 113), confirming a small but
consistent positive I to P relationship globally. H1 (between-regime
Q\_M test) is confirmed (\emph{Q}\_M = 17.35, \emph{p} = .002); however,
the directional Exploratory Propositions E1a (Advanced largest) and E1b
(Frontier smallest) are not statistically confirmed, the Q\_M is driven
by a \emph{k} = 3 Frontier-group anomaly rather than a monotone
institutional gradient. cDAI (\emph{Q}\_M = 1.23, \emph{p} = .541) and
DPL phase (\emph{Q}\_M = 0.62, \emph{p} = .734) are non-significant (H2,
H3 not supported). The most substantive finding is publication bias (H4
confirmed): trim-and-fill imputes \emph{k} = 57 missing studies,
reducing the adjusted estimate to \emph{r} = 0.035, positive but
substantially smaller than the raw pooled effect. The heterogeneity
puzzle (\emph{I}² = 62.4\%) remains unresolved by the tested moderators,
suggesting that future research should test alternative theoretical
contingencies or that between-study heterogeneity is predominantly
within-paper (Level 2, 54.1\%) rather than between-country (Level 3,
8.4\%).

\textbf{Organization.} The paper proceeds as follows. Section 2 develops
the theoretical framework and hypotheses, grounding each moderator in
RBV, institutional theory, organizational learning theory, coordination
cost theory, and the Foundational Digital Adoption Framework. Section 3
describes the systematic search protocol, coding procedure, and
three-level MARA specification. Section 4 presents results, including
the baseline model, three moderator analyses, and publication bias
diagnostics. Section 5 discusses theoretical and practical implications,
and Section 6 concludes with limitations and directions for future
research.

\hypertarget{theoretical-framework-and-hypotheses}{%
\subsection{2. Theoretical Framework and
Hypotheses}\label{theoretical-framework-and-hypotheses}}

\hypertarget{foundation-theories}{%
\subsubsection{2.1 Foundation theories}\label{foundation-theories}}

Five perspectives connect the three new moderators to established I to P
predictions. The \textbf{resource-based view} (Barney, 1991; Wernerfelt,
1984) implies that firms with stronger home-country resource endowments,
including access to national digital infrastructure, earn higher returns
from international expansion; a richer national digital environment
(cDAI) lowers the per-unit cost of deploying firm capabilities abroad.
\textbf{Institutional theory} (North, 1990; Scott, 1995) treats
institutional quality, rule of law, contract enforcement, IP protection,
as the determinant of cross-border coordination costs, yielding the
prediction that the I to P effect declines as institutional quality
falls. \textbf{Organizational learning theory} (Johanson \& Vahlne,
1977, 2009) holds that digital platforms compress the experiential
learning curve once they reach maturity (the DPL Follow phase).
\textbf{Coordination cost theory} (Hitt et al., 1997; Lu \& Beamish,
2004) explains the modal inverted-U; mature national digital
infrastructure attenuates the right-side decline through communication
compression, transaction-cost reduction, and information-asymmetry
compression. Finally, the \textbf{foundational digital adoption
framework} (Verhoef et al., 2021; Stallkamp \& Schotter, 2021)
distinguishes Tier-1/Tier-2 adoption, which functions as shared platform
infrastructure, from proprietary capability; aggregated to the country
level, it defines cDAI.

Throughout, cDAI (country-level) is analytically distinct from
firm-level digital adoption (DAI) and technological capability (TCI),
which the companion primary studies (P3--P8) treat as separate
within-study constructs. At meta-analytic resolution only cDAI is
extractable as a study-level moderator; the claims below are
between-study moderation claims at the country level, not within-firm
capability-bundle claims.

\hypertarget{capabilityinstitution-mismatch-theory-and-the-icrv-gradient-h1}{%
\subsubsection{2.2 Capability--Institution Mismatch Theory and the ICRV
gradient
(H1)}\label{capabilityinstitution-mismatch-theory-and-the-icrv-gradient-h1}}

We propose \emph{Capability--Institution Mismatch Theory} (CIMT): the
productivity return to international expansion depends on the degree to
which home-country institutions let firm capabilities be deployed
productively abroad. Three mechanisms operate. A \emph{rent-protection}
mechanism (Kogut \& Zander, 1993; Zaheer, 1995): strong IP protection
and contract enforcement (ICRV Regime I) let firms retain rents across
markets, raising the average effect. A \emph{liability-of-foreignness
attenuation} mechanism (Peng et al., 2008; Zaheer, 1995): transparent
regulation and low corruption reduce LOF, enlarging captured gains. An
\emph{institutional-void amplification} mechanism (Khanna \& Palepu,
2010): as institutions weaken, substitute-governance costs accumulate,
depressing the effect toward zero. CIMT therefore predicts
between-regime heterogeneity, largest at the institutional frontier.
Prior meta-analyses lacked a taxonomy fine enough to test this; Marano
et al.~(2016) used a coarse six-category scheme. The ICRV 5-regime
taxonomy, anchored in World Bank WGI Rule of Law (2023) and applied to
the full k = 238 corpus (Singapore WGI +1.84 to Pakistan −0.55, Iran
−0.74), provides the first such test.

\textbf{Hypothesis 1 (ICRV between-regime heterogeneity):} Pooled I to P
effects vary systematically across ICRV regimes, with Advanced-regime
studies expected to show the largest average effects (Khanna \& Palepu,
2010; North, 1990; Zaheer, 1995). Formally, the between-regime Q\_M test
is significant (\emph{p} \textless{} .05) and the ICRV-I point estimate
exceeds those for Emerging (ICRV-III) and Mixed (MX).

\emph{Exploratory Proposition E1a:} Advanced-regime (ICRV-I) studies
show the largest pooled effect (direction confirmatory; magnitude not
pre-specified). \emph{Exploratory Proposition E1b:} Frontier-regime
(ICRV-FR) studies show the smallest, potentially null/negative, effect;
treated as exploratory because current Frontier k = 3 is below the k ≥
10 needed for stable moderation (Valentine et al., 2010).

\hypertarget{digital-paradox-lifecycle-h2}{%
\subsubsection{2.3 Digital Paradox Lifecycle
(H2)}\label{digital-paradox-lifecycle-h2}}

The Digital Paradox Lifecycle (DPL) is a temporal moderator grounded in
Brynjolfsson et al.'s (2021) productivity J-curve and David's (1990)
dynamo analogy: general-purpose technologies require years of
complementary investment before productivity gains materialize. Three
phases characterize digital internationalization: \textbf{Precede} (data
predominantly pre-2009), low digital penetration, traditional
coordination-cost dynamics and the inverted-U dominate; \textbf{Span}
(2005--2014), transitional and heterogeneous; \textbf{Follow}
(post-2014), digital platforms reach maturity and compress coordination
costs, so effects should be largest. The 2009 inflection is anchored
empirically in the 2007--2009 smartphone/mobile-internet diffusion, the
2008--2011 rise of Asia-Pacific B2B e-commerce (Alibaba International,
Lazada, Tokopedia), and post-2009 SME digital adoption; the Asia-Pacific
region is 52\% of the corpus.

\textbf{Hypothesis 2 (DPL phase):} Studies on post-2014 data (Follow)
show larger pooled I to P effects than pre-2009 (Precede) studies, with
Span intermediate (Brynjolfsson et al., 2021; David, 1990); the
between-phase Q\_M test is expected significant (\emph{p} \textless{}
.05). Detection is bounded by within-phase k (below \textasciitilde30
per phase the test is underpowered).

\hypertarget{cdai-as-amplifier-h3}{%
\subsubsection{2.4 cDAI as amplifier (H3)}\label{cdai-as-amplifier-h3}}

National digital adoption (cDAI) is an ecosystem property, broadband,
electronic payments, digital identity, regulatory support, that exists
independently of any firm's choice and varies between studies. It is
measured on a 0--1 scale from the World Bank Digital Adoption Index
(2016; Sahay et al., 2020), with the ITU Digital Development Index
(rescaled) as substitute; country-year assignment follows each study's
median data year. The amplification operates through the
coordination-platform effect (Stallkamp \& Schotter, 2021): where
Tier-1/Tier-2 tools are diffused, the per-unit cost of cross-border
transactions falls, especially for higher-intensity internationalizers
(Bharadwaj et al., 2013; Verhoef et al., 2021). Bustamante et al.~(2022)
provide the closest prior cross-country evidence; we extend it
meta-analytically and bound it by DPL phase.

\textbf{Hypothesis 3 (cDAI amplification):} Studies from high-cDAI
contexts show larger pooled I to P effects than low-cDAI studies
(between-group Q\_M significant, \emph{p} \textless{} .05),
operationalized as a Low/Medium/High comparison; the effect is expected
to be concentrated in Follow-phase studies.

\hypertarget{publication-bias-as-a-formal-null-h4}{%
\subsubsection{2.5 Publication bias as a formal null
(H4)}\label{publication-bias-as-a-formal-null-h4}}

Given selective reporting in IB meta-analyses (Borenstein et al., 2021;
Dickersin, 1990), we test publication bias formally. \textbf{Hypothesis
4:} Selective reporting inflates the raw pooled effect. (H4a) Funnel
asymmetry tests (Egger, Begg) are expected significant; (H4b)
trim-and-fill (Duval \& Tweedie, 2000) imputes left-side studies and
yields a smaller but still positive adjusted estimate; (H4c) Orwin's
(1983) fail-safe N exceeds the threshold needed to render the effect
negligible.

\hypertarget{conceptual-model}{%
\subsubsection{2.6 Conceptual model}\label{conceptual-model}}

\emph{Figure 1} presents the three-level MARA with ICRV (H1), cDAI (H3),
and DPL (H2) as study-level moderators of the baseline I to P pooled
effect, with publication bias (H4) assessed at the synthesis level. The
model nests K = 288 effects within k = 238 studies (σ²\_within =
0.00878, I²₍₂₎ = 54.1\%) within between-study heterogeneity (σ²\_between
= 0.00136, I²₍₃₎ = 8.4\%); total I² = 62.4\%.

\hypertarget{method}{%
\subsection{3. Method}\label{method}}

The approach follows the APA Meta-Analysis Reporting Standards (Cooper,
2010) and PRISMA 2020 (Page et al., 2021); the analysis plan was
registered on OSF as a transparency registration (the working corpus had
already been assembled). Three-level MARA was chosen over single-level
random effects because the corpus contains multiple effects per study,
violating independence (Cheung, 2014; Van den Noortgate et al., 2013);
the model decomposes total heterogeneity into within-study (σ²₍₂₎) and
between-study (σ²₍₃₎) components.

\hypertarget{search-strategy}{%
\subsubsection{3.1 Search strategy}\label{search-strategy}}

The primary search covered Web of Science (SSCI, SCI-E, ESCI) and
Scopus, the most comprehensive databases for IB research (Kraus et al.,
2022), supplemented by ABI/INFORM, Business Source Complete,
ScienceDirect, SpringerLink, and Emerald Insight, plus backward and
forward citation tracking from five anchor meta-analyses (Bausch \&
Krist, 2007; Kirca et al., 2012; Marano et al., 2016; Schwens et al.,
2018; Arte \& Larimo, 2022). Topic terms combined internationalization
measures (internationalization/internationalisation, multinationality,
degree of internationalization, international/geographic
diversification, foreign sales/FSTS, foreign assets/FATA, export
intensity, foreign market entry, foreign subsidiaries) with performance
measures (firm/financial/business performance, ROA, ROE, ROS, Tobin's Q,
profitability, labor/total-factor productivity, efficiency) and an
analysis term (correlation, regression, coefficient, effect size). The
Scopus equivalent recalled 29/30 known-item papers (97\%). Temporal
coverage was January 1977--March 2026 (Rugman, 1976; Vernon, 1971).

\hypertarget{eligibility-and-selection}{%
\subsubsection{3.2 Eligibility and
selection}\label{eligibility-and-selection}}

The two authors independently applied criteria in two stages
(title/abstract, then full text), resolving disagreements by discussion.
Included: private-sector firms (government equity ≤ 50\%; financial
sector excluded) with internationalization measured by FSTS, entropy,
foreign-market count, transnationality, or FDI ratio, and performance
measured by accounting (ROA/ROE/ROS), market (Tobin's Q, returns), or
productivity (labor productivity, TFP) indicators, reporting an
extractable effect (r; β convertible via Peterson \& Brown, 2005; t or F
with df). The main analysis was restricted to peer-reviewed journal
articles and in-press articles with DOI; dissertations, theses,
working/conference papers, book chapters, and reports were excluded to
hold peer-review standards constant (grey literature logged in the
PRISMA flow only). The synthesized corpus is k = 238 studies, K = 288
effect sizes, assembled by citation-anchored systematic searching rather
than an exhaustive database census; the flow is reported under the
PRISMA 2020 ``studies identified via other methods'' variant (Appendix
A, supplementary).

\hypertarget{extraction-and-reliability}{%
\subsubsection{3.3 Extraction and
reliability}\label{extraction-and-reliability}}

Parameters (N, focal r or convertible statistic, data-year range,
country/region, DOI and performance operationalization, moderator
features) were extracted using a standardized form. When r was not
reported, conversions were applied in precision order: r =
√{[}t²/(t²+df){]} (Cohen, 1988); r\_partial = β × 0.98 (the simplified
form of Peterson \& Brown's (2005) estimator 0.98β + 0.05λ; identical
for negative β and conservative for positive β); r = √{[}F/(F+df₂){]}
with df₁ = 1 (Rosenthal, 1994). For multiple estimates per study,
distinct subsamples were coded as independent effects, and for multiple
specifications on the same sample the most fully controlled model was
retained. Inter-coder reliability used a three-stage design (calibration
on 10 studies; independent coding of a 20\% stratified subsample; PI
adjudication), with Cohen's κ (target ≥ 0.70; Landis \& Koch, 1977) for
categorical moderators and ICC(2,1) (≥ 0.80) for continuous cDAI.
Consistent with prior I to P meta-analyses (Bausch \& Krist, 2007;
Marano et al., 2016), no study-level risk-of-bias instrument was
applied; the salient threats, publication bias, omitted-variable bias,
and measurement heterogeneity, are addressed at the synthesis level.

\hypertarget{moderators}{%
\subsubsection{3.4 Moderators}\label{moderators}}

Seven moderators were coded: four standard (country of origin; industry
sector; DOI operationalization; performance operationalization) and
three novel. \emph{ICRV regime} uses World Bank WGI Rule of Law (2023):
I Advanced-Innovation (WGI \textgreater{} +0.80), II Upper-Middle
(0--+0.80), III Emerging (−0.50--0), FR Frontier/SIDS (≤ −0.50 or small
island state), MX multi-country pooled (no modal regime ≥ 60\%).
\emph{cDAI} is the 0--1 World Bank Digital Adoption Index (2016) or
rescaled ITU index, by median data year. \emph{DPL phase} is Precede
(median year \textless{} 2009), Span (2005--2014 or unclassifiable), or
Follow (median year ≥ 2015).

\hypertarget{three-level-mara}{%
\subsubsection{3.5 Three-level MARA}\label{three-level-mara}}

The model decomposes effect r\_ij (effect i in study j) into sampling
error, within-study, and between-study components:

\[r_{ij} = \theta_{ij} + e_{ij}, \quad e_{ij} \sim \mathcal{N}(0, v_{ij}), \quad v_{ij} \approx 1/(N_{ij}-3)\]
\[\theta_{ij} = \delta_{j} + u_{ij}, \quad u_{ij} \sim \mathcal{N}(0, \sigma_{(2)}^{2})\]
\[\delta_{j} = \mu + \mathbf{X}_{j}\mathbf{\beta} + w_{j}, \quad w_{j} \sim \mathcal{N}(0, \sigma_{(3)}^{2})\]

where \textbf{X}\_j holds the study-level moderators (ICRV dummies,
cDAI, DPL dummies) plus the four standard controls. Parameters are
estimated by REML via \texttt{rma.mv} in \texttt{metafor} v4
(Viechtbauer, 2010) with a compound-symmetric variance structure for
within-study effects; REML is preferred for unbiased variance components
at moderate k (Raudenbush \& Bryk, 2002). All r are Fisher-z transformed
before analysis and back-transformed for reporting (Hedges \& Olkin,
1985). Level-specific I² is computed per Cheung (2014). Categorical
moderators use the Q\_M omnibus test (Holm--Bonferroni for pairwise
regime comparisons); continuous cDAI uses a two-sided Wald z-test.

\hypertarget{publication-bias-and-robustness}{%
\subsubsection{3.6 Publication bias and
robustness}\label{publication-bias-and-robustness}}

Publication bias was assessed with Egger's regression (Egger et al.,
1997), Begg and Mazumdar's (1994) rank correlation, Duval and Tweedie's
(2000) trim-and-fill, Orwin's (1983) fail-safe N, and PET-PEESE
meta-regression (Stanley \& Doucouliagos, 2014). Five pre-specified
robustness checks follow: (1) two-level vs.~three-level comparison (Δr
\textgreater{} 0.02 flags nesting bias); (2) leave-one-out influence
(Cook's distance \textgreater{} 4/k); (3) FSTS-only restriction; (4)
ICRV reclassified on the WGI composite governance index; (5) post-2000
temporal restriction to test vintage effects.

\hypertarget{results}{%
\subsection{4. Results}\label{results}}

\hypertarget{sample-description-and-inter-coder-reliability}{%
\subsubsection{4.1 Sample Description and Inter-Coder
Reliability}\label{sample-description-and-inter-coder-reliability}}

\emph{k} = 238 studies (coded), \emph{K} = 288 effect sizes.

\textbf{Table 3.1, Inter-coder reliability} \emph{(two authors
independently code a 20\% stratified subsample, k = 47 studies, blind to
each other's codes)}

\begin{longtable}[]{@{}
  >{\raggedright\arraybackslash}p{(\columnwidth - 8\tabcolsep) * \real{0.2800}}
  >{\raggedright\arraybackslash}p{(\columnwidth - 8\tabcolsep) * \real{0.2400}}
  >{\raggedright\arraybackslash}p{(\columnwidth - 8\tabcolsep) * \real{0.1467}}
  >{\raggedright\arraybackslash}p{(\columnwidth - 8\tabcolsep) * \real{0.0933}}
  >{\raggedright\arraybackslash}p{(\columnwidth - 8\tabcolsep) * \real{0.2400}}@{}}
\toprule\noalign{}
\begin{minipage}[b]{\linewidth}\raggedright
Moderator
\end{minipage} & \begin{minipage}[b]{\linewidth}\raggedright
Variable type
\end{minipage} & \begin{minipage}[b]{\linewidth}\raggedright
Statistic
\end{minipage} & \begin{minipage}[b]{\linewidth}\raggedright
Value
\end{minipage} & \begin{minipage}[b]{\linewidth}\raggedright
Target threshold
\end{minipage} \\
\midrule\noalign{}
\endhead
\bottomrule\noalign{}
\endlastfoot
ICRV regime & Categorical (5) & Cohen's κ & {[}insert after
dual-coding{]} & ≥ 0.70 \\
DPL phase & Categorical (3) & Cohen's κ & {[}insert after dual-coding{]}
& ≥ 0.70 \\
Industry sector & Categorical (3) & Cohen's κ & {[}insert after
dual-coding{]} & ≥ 0.70 \\
DOI measure & Categorical (4) & Cohen's κ & {[}insert after
dual-coding{]} & ≥ 0.70 \\
Performance measure & Categorical (4) & Cohen's κ & {[}insert after
dual-coding{]} & ≥ 0.70 \\
cDAI score & Continuous (0--1) & ICC(2, 1) & {[}insert after
dual-coding{]} & ≥ 0.80 \\
\end{longtable}

\emph{Note.} The full corpus is coded by the first author against the
standardized codebook (Appendix B). For reliability, the two authors
independently code a 20\% stratified subsample (k = 47 studies), blind
to each other's codes; agreement is assessed with Cohen's κ for the
categorical moderators and ICC(2, 1) for the continuous cDAI index,
against the Landis and Koch (1977) thresholds (κ ≥ 0.70; ICC ≥ 0.80).
Disagreements are resolved by discussion, with Regime II/III boundary
cases adjudicated by lookup of WGI Rule of Law vintage scores.

\textbf{Table 4.1, Sample composition} \emph{(K = 288 effect sizes, k =
238 coded studies)}

\begin{longtable}[]{@{}
  >{\raggedright\arraybackslash}p{(\columnwidth - 4\tabcolsep) * \real{0.7374}}
  >{\raggedright\arraybackslash}p{(\columnwidth - 4\tabcolsep) * \real{0.1313}}
  >{\raggedright\arraybackslash}p{(\columnwidth - 4\tabcolsep) * \real{0.1313}}@{}}
\toprule\noalign{}
\begin{minipage}[b]{\linewidth}\raggedright
Category
\end{minipage} & \begin{minipage}[b]{\linewidth}\raggedright
\emph{K} effects
\end{minipage} & \begin{minipage}[b]{\linewidth}\raggedright
\emph{k} studies
\end{minipage} \\
\midrule\noalign{}
\endhead
\bottomrule\noalign{}
\endlastfoot
ICRV Regime I, Advanced (e.g., Korea, Japan, Singapore, HK, Australia) &
139 & 107 \\
ICRV Regime II, Upper-middle (e.g., China, Malaysia, Thailand) & 25 &
21 \\
ICRV Regime III, Emerging (e.g., Vietnam, India, Philippines) & 91 &
79 \\
ICRV Frontier / SIDS (FR) & 3 & 3 \\
Cross-regime / multi-country (MX) & 30 & 28 \\
\textbf{\emph{K} / \emph{k} total} & \textbf{288} & \textbf{238} \\
cDAI High (H) & 38 & , \\
cDAI Medium (M) & 76 & , \\
cDAI Low (L) & 174 & , \\
DPL Precede (PRE, ≤2008) & 100 & , \\
DPL Span (SPN, 2009--2013) & 108 & , \\
DPL Follow (FOL, ≥2014) & 80 & , \\
By DOI type: FSTS & 138 & , \\
By DOI type: GEO & 50 & , \\
By DOI type: EXP & 65 & , \\
By DOI type: COMP & 31 & , \\
By FP type: ACC (accounting) & 246 & , \\
By FP type: MKT (market-based) & 16 & , \\
By FP type: LAB (labour productivity) & 12 & , \\
By FP type: MIX & 14 & , \\
\end{longtable}

\emph{Note:} Counts from the coded database
(\texttt{p6/results/forest\_data.csv}, K=288 rows, k=238 unique study
IDs, updated 15/05/2026). ICRV \emph{k} and \emph{K} counts sum to
\textgreater{} total because MX studies may span multiple regimes. Study
(\emph{k}) counts by cDAI/DPL are reported after multi-effect
deduplication. Within the combined Frontier/SIDS (FR) code, all \emph{k} =
3 coded studies sample frontier economies and none sample a small-island
developing state (SIDS \emph{k} = 0); SIDS effects therefore cannot be
estimated in the present corpus and are flagged as a targeted-search
priority (Appendix B).

\hypertarget{baseline-three-level-model}{%
\subsubsection{4.2 Baseline Three-Level
Model}\label{baseline-three-level-model}}

\textbf{ICBEF 2025 single-level baseline (MetaEssentials 1.5, k = 113):}

\[{\bar{r}}_{ICBEF} = 0.07\quad\left( 95\%\ \text{CI}:\lbrack 0.05, 0.09\rbrack \right), \ p < .001\]

\[I_{ICBEF}^{2} = 87.92\%, \quad Q_{between} = 1,247.3\ (df = 112, \ p < .001)\]

\textbf{Three-level decomposition} (metafor REML, k = 238 studies, K =
288 effects):

\begin{longtable}[]{@{}ll@{}}
\toprule\noalign{}
Parameter & Estimate \\
\midrule\noalign{}
\endhead
\bottomrule\noalign{}
\endlastfoot
σ²\_(2) within-study & 0.00878 \\
σ²\_(3) between-study & 0.00136 \\
\emph{I}²\_(2) within-study & 54.1\% \\
\emph{I}²\_(3) between-study & 8.4\% \\
\emph{I}²\_total & 62.4\% \\
Pooled \emph{r̂}\_3L & 0.074 (95\% CI {[}0.060, 0.088{]}) \\
\emph{Q}\_total & 1,895.58 (\emph{df} = 286, \emph{p} \textless{}
.001) \\
\end{longtable}

The three-level pooled estimate (\emph{r̂} = 0.074) is consistent with
the ICBEF 2025 single-level baseline (\emph{r} = 0.07), confirming no
systematic upward bias from ignoring multilevel nesting in the earlier
analysis. Three-level decomposition reveals that the larger share of
systematic heterogeneity lies within studies (Level 2, \emph{I}²\_(2) =
54.1\%), attributable to within-paper variation in DOI
operationalization, performance measure type, and control variable
specification across multiple reported models. Between-study variance
(Level 3, \emph{I}²\_(3) = 8.4\%) represents country-context differences
that ICRV, cDAI, and DPL phase were designed to explain. Total \emph{I}²
= 62.4\%, indicating substantial heterogeneity beyond sampling error,
the motivation for moderator analysis, though the moderator analyses in
Section 4.3--4.5 do not resolve this heterogeneity significantly.

\hypertarget{icrv-5-regime-moderation-h1}{%
\subsubsection{4.3 ICRV 5-Regime Moderation
(H1)}\label{icrv-5-regime-moderation-h1}}

\emph{Q}\_M(ICRV) = 17.35 (\emph{df} = 4, \emph{p} = .002). H1
\textbf{confirmed}, the between-regime Q\_M test is statistically
significant as hypothesized. Exploratory Propositions E1a (Advanced
largest) and E1b (Frontier smallest) are \textbf{not meta-analytically
confirmed}: the significant Q\_M is driven by the \emph{k} = 3
Frontier-group anomaly rather than a monotone institutional gradient;
see below.

\textbf{Table 4.1, ICRV regime subgroup results} \emph{(actual MARA
output, k=238/K=288)}

\begin{longtable}[]{@{}
  >{\raggedright\arraybackslash}p{(\columnwidth - 8\tabcolsep) * \real{0.2881}}
  >{\raggedright\arraybackslash}p{(\columnwidth - 8\tabcolsep) * \real{0.0282}}
  >{\raggedright\arraybackslash}p{(\columnwidth - 8\tabcolsep) * \real{0.0395}}
  >{\raggedright\arraybackslash}p{(\columnwidth - 8\tabcolsep) * \real{0.1017}}
  >{\raggedright\arraybackslash}p{(\columnwidth - 8\tabcolsep) * \real{0.5424}}@{}}
\toprule\noalign{}
\begin{minipage}[b]{\linewidth}\raggedright
Regime
\end{minipage} & \begin{minipage}[b]{\linewidth}\raggedright
\emph{k}
\end{minipage} & \begin{minipage}[b]{\linewidth}\raggedright
\emph{r̄}
\end{minipage} & \begin{minipage}[b]{\linewidth}\raggedright
95\% CI
\end{minipage} & \begin{minipage}[b]{\linewidth}\raggedright
Note
\end{minipage} \\
\midrule\noalign{}
\endhead
\bottomrule\noalign{}
\endlastfoot
I, Advanced-Innovation (SG, HK, KR, JP, UK, US\ldots) & 139 & 0.079 &
{[}0.058, 0.099{]} & Largest group; broadly consistent with H1 \\
II, Upper-middle & 25 & 0.065 & {[}0.020, 0.109{]} & ns vs.~Group I (b =
−0.014, \emph{p} = .581) \\
III, Emerging (VN, IN, CN, PH\ldots) & 91 & 0.068 & {[}0.044, 0.092{]} &
ns vs.~Group I (b = −0.011, \emph{p} = .502) \\
FR, Frontier & 3 & 0.349 & {[}0.217, 0.468{]} & ⚠️ \emph{k} = 3 only;
estimate unreliable (one outlier: Pouresmaeili et al., \emph{r} = 0.69,
\emph{n} = 226) \\
MX, Multi-country / Mixed & 30 & 0.053 & {[}0.012, 0.094{]} & ns
vs.~Group I (b = −0.026, \emph{p} = .269) \\
\end{longtable}

Figure 2: ICRV 5-regime forest plot, subgroup pooled effects with 95\%
CI

\emph{Figure 2.} ICRV subgroup forest plot. Pooled \emph{r̄} and 95\% CI
per coded institutional regime.

The \emph{Q}\_M = 17.35 is statistically significant (\emph{p} = .002)
but the between-group pattern does not confirm the hypothesized monotone
gradient (I \textgreater{} II \textgreater{} III \textgreater{}
Frontier). Pairwise contrasts against Group I (intercept) reveal that
Group II (\emph{p} = .581), Group III (\emph{p} = .502), and MX
(\emph{p} = .269) do not differ significantly from Advanced-Innovation
contexts. The only significant contrast is Group I vs.~Frontier (b =
+0.285, \emph{p} \textless{} .001), driven entirely by the \emph{k} = 3
Frontier-group estimate, which is dominated by one outlier study
(Pouresmaeili et al., 2018, \emph{r} = 0.69). Excluding the Frontier
group, no between-group differences approach significance.

Group I (k=139, primarily Western MNEs and Asian advanced economies)
shows the highest reliable I to P effect (\emph{r̄} = 0.079), and Group
III (k=91, Emerging) shows \emph{r̄} = 0.068, a difference of 0.011 that
is directionally consistent with CIMT but not statistically significant
given within-group heterogeneity. H1 (between-regime heterogeneity
\emph{Q}\_M significant) is confirmed, establishing that I to P effect
sizes vary systematically across ICRV regimes. However, E1a's predicted
directional ordering (Advanced \textgreater{} Emerging) and E1b's
predicted Frontier floor are not statistically confirmable at \emph{k} =
3 Frontier. The gradient-specific propositions require larger \emph{k}
per regime before they can be properly tested.

\hypertarget{cdai-moderation-h3}{%
\subsubsection{4.4 cDAI Moderation (H3)}\label{cdai-moderation-h3}}

\emph{Q}\_M(cDAI) = 1.23 (\emph{df} = 2, \emph{p} = .541). H3
\textbf{not supported}.

Meta-regression with continuous cDAI score: β\_cDAI = +0.003 (\emph{SE}
= 0.010, \emph{p} = .744). Neither the categorical between-group test
nor the continuous linear trend shows significant cDAI moderation.

\textbf{Table 4.2, cDAI subgroup} \emph{(actual MARA output)}

\begin{longtable}[]{@{}lllll@{}}
\toprule\noalign{}
cDAI Group & \emph{k} & \emph{r̄} & 95\% CI & Δ vs.~Low \\
\midrule\noalign{}
\endhead
\bottomrule\noalign{}
\endlastfoot
Low & 174 & 0.075 & {[}0.056, 0.094{]} & , \\
Medium & 76 & 0.063 & {[}0.036, 0.090{]} & b = −0.012, \emph{p} =
.489 \\
High & 38 & 0.091 & {[}0.052, 0.129{]} & b = +0.016, \emph{p} = .469 \\
\end{longtable}

The three cDAI subgroup means are all significantly positive (\emph{p}
\textless{} .001) but do not differ significantly from each other. The
non-monotone ordering (Low \textgreater{} Medium \textless{} High) and
small, non-significant contrasts fail to support the predicted positive
linear gradient. The cDAI × DPL interaction cannot be reliably estimated
in the current sample given the non-significant main moderation. H3 is
not supported: country-level digital adoption (cDAI, measured via World
Bank DAI / ITU Digital Development Index) does not significantly amplify
the pooled I to P effect in this \emph{k} = 238 sample. A larger sample
with better resolution across the cDAI spectrum is required to test the
cDAI amplification (H3) gradient hypothesis.

\hypertarget{dpl-phase-moderation-h2}{%
\subsubsection{4.5 DPL Phase Moderation
(H2)}\label{dpl-phase-moderation-h2}}

\emph{Q}\_M(DPL) = 0.62 (\emph{df} = 2, \emph{p} = .734). H2 \textbf{not
supported}.

\textbf{Table 4.3, DPL phase subgroup} \emph{(actual MARA output)}

\begin{longtable}[]{@{}
  >{\raggedright\arraybackslash}p{(\columnwidth - 10\tabcolsep) * \real{0.1415}}
  >{\raggedright\arraybackslash}p{(\columnwidth - 10\tabcolsep) * \real{0.3491}}
  >{\raggedright\arraybackslash}p{(\columnwidth - 10\tabcolsep) * \real{0.0472}}
  >{\raggedright\arraybackslash}p{(\columnwidth - 10\tabcolsep) * \real{0.0660}}
  >{\raggedright\arraybackslash}p{(\columnwidth - 10\tabcolsep) * \real{0.1698}}
  >{\raggedright\arraybackslash}p{(\columnwidth - 10\tabcolsep) * \real{0.2264}}@{}}
\toprule\noalign{}
\begin{minipage}[b]{\linewidth}\raggedright
DPL Phase
\end{minipage} & \begin{minipage}[b]{\linewidth}\raggedright
Definition
\end{minipage} & \begin{minipage}[b]{\linewidth}\raggedright
\emph{k}
\end{minipage} & \begin{minipage}[b]{\linewidth}\raggedright
\emph{r̄}
\end{minipage} & \begin{minipage}[b]{\linewidth}\raggedright
95\% CI
\end{minipage} & \begin{minipage}[b]{\linewidth}\raggedright
Δ vs.~PRE
\end{minipage} \\
\midrule\noalign{}
\endhead
\bottomrule\noalign{}
\endlastfoot
Precede (PRE) & Sample data predominantly pre-2009 & 100 & 0.082 &
{[}0.057, 0.107{]} & , \\
Span (SPN) & Sample spans the 2009 inflection & 108 & 0.068 & {[}0.045,
0.091{]} & b = −0.014, \emph{p} = .434 \\
Follow (FOL) & Sample data predominantly post-2014 & 80 & 0.073 &
{[}0.046, 0.100{]} & b = −0.009, \emph{p} = .645 \\
\end{longtable}

Pairwise comparisons: PRE vs.~FOL (\emph{z} = 0.46, \emph{p} = .645);
PRE vs.~SPN (\emph{z} = 0.78, \emph{p} = .434); FOL vs.~SPN (\emph{z} =
0.28, \emph{p} = .782). No pairwise difference approaches significance.

The ordering (PRE \textgreater{} FOL \textgreater{} SPN) is opposite to
H2's prediction (FOL \textgreater{} SPN \textgreater{} PRE). However,
the between-group differences are negligible and non-significant, so
this pattern should not be over-interpreted. The null DPL result may
reflect that the \emph{k} = 238 sample does not have sufficient power to
detect small temporal trends, that DPL phase is confounded with ICRV
composition (pre-2009 studies concentrated in advanced economies), or
that the Digital Paradox Lifecycle inflection does not manifest at
meta-analytic resolution with the current sample size. H2 is not
supported.

Figure 3: DPL phase moderation, pooled I to P effect by digital adoption
epoch

\emph{Figure 3.} DPL phase subgroup results. Pooled effect sizes by
Precede / Span / Follow epoch with 95\% CI. The between-phase
differences are small and non-significant.

\hypertarget{publication-bias-h4}{%
\subsubsection{4.6 Publication Bias (H4)}\label{publication-bias-h4}}

H4 \textbf{supported}: multiple indicators consistently detect
publication bias, though the positive I to P effect survives correction.

\textbf{Egger's regression test} (SE as moderator): \emph{b} = 0.487
(\emph{SE} = 0.251, \emph{p} = .052), just above the conventional α =
.05 threshold; this criterion does not reach statistical significance,
and funnel asymmetry by Egger's test alone is not confirmed.

\textbf{Begg's rank correlation} (Kendall's τ): τ = −0.132, \emph{p} =
.001 significant funnel asymmetry; studies with larger standard errors
(smaller \emph{n}) report systematically lower effect sizes, consistent
with publication bias against null or negative findings.

\textbf{Trim-and-fill}: imputes \emph{k} = 57 missing studies (left
side); adjusted pooled \emph{r̄} = 0.035 (95\% CI {[}0.018, 0.051{]}), a
conservative lower bound. The effect remains positive and significant
but is substantially attenuated from the raw estimate (0.074 to 0.035).

\textbf{Fail-safe \emph{N}} (Rosenthal, 1991): \emph{N} = 44,782, far
exceeding the criterion of 5\emph{k} + 10 = 1,195; even under extreme
publication suppression assumptions, a trivially small effect would
require 44,782 unpublished null studies, implausible.

The trim-and-fill correction (\emph{k} = 57 imputed, adj. \emph{r} =
0.035) is the most conservative bias-corrected estimate and represents a
meaningful reduction from the raw \emph{r} = 0.074. Together with the
substantial unexplained heterogeneity (\emph{I}² = 62.4\%) and
non-significant moderator tests (Section 4.3--4.5), the publication bias
evidence suggests that the apparent average I to P effect is upwardly
inflated in the published literature. The true population effect may be
closer to \emph{r} ≈ 0.035.

Figure 5: Funnel plot with trim-and-fill imputed studies

\emph{Figure 5.} Funnel plot of effect sizes against standard errors.
Open circles = original studies; filled circles = trim-and-fill imputed
studies (\emph{k} = 57). Substantial left-side asymmetry is visible; the
adjusted pooled effect is \emph{r} = 0.035.

\hypertarget{robustness}{%
\subsubsection{4.7 Robustness}\label{robustness}}

\begin{longtable}[]{@{}
  >{\raggedright\arraybackslash}p{(\columnwidth - 8\tabcolsep) * \real{0.3604}}
  >{\raggedright\arraybackslash}p{(\columnwidth - 8\tabcolsep) * \real{0.0450}}
  >{\raggedright\arraybackslash}p{(\columnwidth - 8\tabcolsep) * \real{0.1622}}
  >{\raggedright\arraybackslash}p{(\columnwidth - 8\tabcolsep) * \real{0.1622}}
  >{\raggedright\arraybackslash}p{(\columnwidth - 8\tabcolsep) * \real{0.2703}}@{}}
\toprule\noalign{}
\begin{minipage}[b]{\linewidth}\raggedright
Check
\end{minipage} & \begin{minipage}[b]{\linewidth}\raggedright
K
\end{minipage} & \begin{minipage}[b]{\linewidth}\raggedright
\emph{r̄}
\end{minipage} & \begin{minipage}[b]{\linewidth}\raggedright
95\% CI
\end{minipage} & \begin{minipage}[b]{\linewidth}\raggedright
Note
\end{minipage} \\
\midrule\noalign{}
\endhead
\bottomrule\noalign{}
\endlastfoot
Main analysis & 288 & 0.074 & {[}0.060, 0.088{]} & Baseline \\
Confirmed \emph{r} only (exclude estimated) & 240 & 0.077 & {[}0.059,
0.094{]} & Consistent \\
Exclude \emph{n} \textless{} 30 & 285 & 0.073 & {[}0.059, 0.088{]} &
Consistent \\
ACC performance only & 246 & 0.075 & {[}0.059, 0.091{]} & Consistent \\
FSTS DOI measure only & 137 & 0.060 & {[}0.041, 0.078{]} & Attenuated
but positive \\
DL estimator (two-level) & 288 & 0.074 & {[}0.061, 0.086{]} & Δ\emph{r}
\textless{} 0.01 vs.~three-level \\
Leave-one-out range & 288 & {[}0.071, 0.075{]} & , & 0/288 change
direction \\
\end{longtable}

The baseline \emph{r} = 0.074 is robust across all sensitivity checks.
The leave-one-out range {[}0.071, 0.075{]} confirms no single study
drives the result. The DL two-level estimator gives an identical point
estimate (0.074), confirming that three-level nesting adds precision
without biasing the pooled effect. The FSTS-only restriction (\emph{r̄} =
0.060) is the most conservative check and remains positive and
significant, suggesting that DOI operationalization heterogeneity
partially inflates the pooled effect when broader DOI measures are
included.

Figure 4: Leave-one-out sensitivity, pooled \emph{r̄} range across
leave-one-out iterations

\emph{Figure 4.} Leave-one-out sensitivity analysis. Each point is the
pooled \emph{r̄} with 95\% CI after removing one study. The narrow range
confirms no single study drives the results.

\hypertarget{discussion}{%
\subsection{5. Discussion}\label{discussion}}

\hypertarget{alignment-with-country-level-evidence}{%
\subsubsection{5.1 Alignment with Country-Level
Evidence}\label{alignment-with-country-level-evidence}}

The meta-analytic baseline (\emph{r} = 0.074, \emph{k} = 238, 49
economies) is consistent with country-level evidence from across the
global study corpus, including the Asia-Pacific primary studies
underlying the ICBEF 2025 baseline (Authors, 2024). The pooled effect is
positive and significant, confirming that export-intensive firms tend to
outperform domestically focused peers even after adjusting for firm
size, age, and industry.

\textbf{Advanced-I contexts (Group I, k = 139; \emph{r̄} = 0.079):} The
Group I mean (\emph{r̄} = 0.079) is the highest reliable estimate and is
directionally consistent with institutional complementarity theory:
strong contract enforcement, IP protection, and low liability of
foreignness (Zaheer, 1995) allow firms to sustain productive
internationalization without the coordination cost penalties visible in
weaker environments (North, 1990; Peng et al., 2008). However, the
magnitude is much smaller than CIMT predicted, and the I--III difference
is not statistically significant (\emph{p} = .502).

\textbf{Emerging contexts (Group III, k = 91; \emph{r̄} = 0.068):} The
Emerging group shows \emph{r̄} = 0.068, directionally below the Advanced
group, consistent with CIMT's institutional friction argument. At the
primary-study level, Authors (2025) document a negative aggregate DOI
coefficient on 380 Indian WBES firms (FGLS estimation), with manager
experience and female top-management leadership as positive moderating
factors, a pattern consistent with CIMT's prediction that firm-level
compensating resources partially offset institutional friction in
Emerging environments.

\textbf{Frontier contexts (FR, k = 3; \emph{r̄} = 0.349):} The \emph{k} =
3 Frontier estimate is driven by one outlier study (Pouresmaeili et al.,
2018, \emph{r} = 0.69, \emph{n} = 226 Iranian manufacturing firms). This
estimate cannot be interpreted as a reliable Frontier I to P effect; it
instead reflects the near-impossibility of drawing meta-analytic
conclusions when only three studies populate a regime cell. The actual
Frontier I to P effect could be zero, positive, or negative, larger
\emph{k} is required.

\textbf{Synthesis:} H1 is confirmed, the ICRV between-regime Q\_M test
(\emph{Q}\_M = 17.35, \emph{p} = .002) is statistically significant,
establishing that I to P effect sizes vary across institutional regimes
as theorized. The Q\_M significance is driven by the \emph{k} = 3
Frontier contrast rather than by well-powered pairwise contrasts across
the Advanced/Upper-middle/Emerging cells, so the directional gradient
propositions E1a and E1b are not meta-analytically confirmed in the
current sample. The substantiated finding is that the baseline I to P
effect (\emph{r} = 0.074) is broadly consistent across well-populated
regime groups (I, II, III, MX), and the magnitude of the
Advanced-Emerging gap (\emph{r̄} = 0.079 vs.~0.068) does not reach
significance, a result that bounds, rather than refutes, CIMT's
institutional gradient prediction.

\hypertarget{theoretical-contributions}{%
\subsubsection{5.2 Theoretical
Contributions}\label{theoretical-contributions}}

\textbf{Contribution 1: Largest globally representative I to P
meta-analysis to date.} The three-level MARA with \emph{k} = 238 studies
from 49 economies (\emph{K} = 288 effects) is the most comprehensive
quantitative synthesis of I to P evidence to apply three-level modeling,
extending the ICBEF 2025 Asia-Pacific baseline (\emph{k} = 113) by 124
studies with a globally inclusive ICRV-coded corpus. Three-level
modeling correctly partitions within-study and between-study variance,
and the baseline \emph{r} = 0.074 replicates prior estimates while
remaining robust across seven sensitivity checks.

\textbf{Contribution 2: First meta-analytic test of ICRV, cDAI, and DPL
moderators.} This paper is the first to formally test ICRV institutional
regime, national digital adoption (cDAI), and Digital Paradox Lifecycle
phase as moderators in a three-level MARA. H1 (between-regime Q\_M test)
is confirmed (\emph{Q}\_M = 17.35, \emph{p} = .002), establishing that I
to P effect sizes vary systematically across institutional regimes.
However, the directional gradient propositions (E1a/E1b) and the cDAI
and DPL hypotheses (H2, H3) are not confirmed under rigorous
between-group testing. These results set theoretically informative
bounds: the \emph{k} = 238 sample lacks sufficient between-regime
variation (particularly Frontier, \emph{k} = 3; SIDS, \emph{k} = 0) to
detect the CIMT gradient, and cDAI and DPL phase do not explain
between-study heterogeneity in the present corpus. Future replications
with targeted regime-level sampling should test whether the gradient
emerges with richer between-regime representation.

\textbf{Contribution 3: Substantial publication bias identified.} The
trim-and-fill-corrected estimate (\emph{r} = 0.035, \emph{k} = 57
imputed studies) and significant Begg's τ (\emph{p} = .001) together
indicate that the I to P literature has a substantial positive
publication bias. The raw pooled effect (\emph{r} = 0.074) likely
overstates the true average causal effect by approximately 50--100\%,
with the bias-corrected estimate (\emph{r} = 0.035) providing a
conservative lower bound. This finding, from the most globally
comprehensive three-level MARA of I to P to date, is the most actionable
theoretical contribution: it suggests that the IB literature's widely
cited ``positive I to P relationship'' is partially an artifact of
selective publication rather than a robust phenomenon.

\hypertarget{managerial-and-policy-implications}{%
\subsubsection{5.3 Managerial and Policy
Implications}\label{managerial-and-policy-implications}}

The non-confirmation of the three hypothesized moderators limits strong
prescriptive conclusions about institutional regime or digital context,
but the baseline finding and publication bias result carry practical
significance.

\textbf{For firms across all institutional contexts:} The bias-corrected
baseline (\emph{r} = 0.035) rather than the raw \emph{r} = 0.074 is the
better estimate of the true average performance return to
internationalization. Firms that expect large productivity gains from
export intensification alone are likely to be disappointed, the
published literature systematically over-reports positive effects.
Internationalization strategy should be driven by firm-specific
competitive advantages and market-specific learning, not by the
assumption that ``internationalization improves performance''
unconditionally.

\textbf{For researchers:} The substantial publication bias (\emph{k} =
57 imputed, adj. \emph{r} = 0.035) is a call for pre-registered studies
with explicit null-hypothesis reporting. Publication practices in the I
to P literature, where positive effects are over-represented, may be
distorting the field's understanding of when and why
internationalization helps. OSF registration and adoption of three-level
MARA with trim-and-fill correction should become standard in future I to
P meta-analyses.

\textbf{For policymakers:} The ICRV Advanced--Emerging gap (\emph{r̄} =
0.079 vs. 0.068, non-significant) is not large enough to justify
institutional quality as the primary determinant of export-led
productivity. Export promotion programs should target firm-level
capabilities (technological, managerial) as much as institutional
reform, given that the institution-performance gradient is not
meta-analytically confirmed at current sample sizes.

\hypertarget{limitations-and-inferential-bounds}{%
\subsubsection{5.4 Limitations and Inferential
Bounds}\label{limitations-and-inferential-bounds}}

Three limitations bound the inferences available from this study:

\textbf{(a) What cannot be concluded:} (1) The SIDS subgroup (\emph{k} =
0; no primary small-island study met inclusion) does not permit any test
of the ``forced-penalty'' hypothesis, a targeted search for SIDS-focused
primary studies (as specified in Appendix B) is required before any
meta-analytic SIDS effect can be estimated. (2) The cDAI × ICRV joint
moderation (three-way interaction) is underpowered in the current
dataset (\emph{k} per cell \textless{} 20); point estimates are provided
but confidence intervals are wide. (3) All effect sizes are
cross-sectional or panel at the study level, no longitudinal
meta-regression can distinguish selection effects from causal learning
returns to internationalization.

\textbf{(b) Methodological remedies for future work:} Panel
meta-analysis with longitudinal effect sizes from individual-firm data
would enable causal decomposition (Sutton \& Higgins, 2008). Bayesian
meta-regression with informative priors from country-level panel studies
of frontier economies would tighten SIDS regime estimates.

\textbf{(c) Boundary of the ICRV classification:} The 5-regime taxonomy
uses WGI Rule of Law as the primary classifier. Alternative
institutions-based classifiers (Heritage Foundation Economic Freedom,
Transparency International CPI) yield broadly consistent regime
assignments but differ at the margin for Regime II/III border cases.

\hypertarget{conclusion}{%
\subsection{6. Conclusion}\label{conclusion}}

This paper presents the first three-level MARA of the
internationalization--performance relationship drawing on a globally
representative corpus, extending the ICBEF 2025 Asia-Pacific baseline
from \emph{k} = 113 to \emph{k} = 238 studies from 49 economies
(\emph{K} = 288 effects) and introducing three novel moderators for
formal meta-analytic testing: ICRV institutional regime, country-level
digital adoption (cDAI), and Digital Paradox Lifecycle (DPL) phase. The
three-level pooled effect (\emph{r} = 0.074, 95\% CI {[}0.060, 0.088{]},
\emph{I}²\_total = 62.4\%) confirms a small but consistent positive I to
P relationship globally, robust across seven sensitivity checks.

H1 (ICRV between-regime Q\_M test) is confirmed (\emph{Q}\_M = 17.35,
\emph{p} = .002), establishing that I to P effect sizes vary
systematically across institutional regimes. However, the directional
gradient (E1a: Advanced largest; E1b: Frontier smallest) is not
meta-analytically confirmable at \emph{k} = 3 Frontier, and the cDAI
(H3: \emph{Q}\_M = 1.23, \emph{p} = .541) and DPL phase (H2: \emph{Q}\_M
= 0.62, \emph{p} = .734) moderators do not explain between-study
heterogeneity in the current sample. The CIMT institutional gradient
(E1a/E1b), digital infrastructure amplification (H3), and DPL temporal
moderation (H2) remain theoretically motivated but empirically
unconfirmed at \emph{k} = 238, null results that should inform future
pre-registered replications with targeted regime-level sampling.

The most substantive empirical finding is substantial publication bias:
trim-and-fill imputes \emph{k} = 57 missing studies, reducing the
bias-corrected pooled effect to \emph{r} = 0.035 (95\% CI {[}0.018,
0.051{]}) positive but considerably attenuated. This suggests that the I
to P field's published literature over-represents positive results, and
that the true average performance return to internationalization is
closer to \emph{r} ≈ 0.035 than the oft-cited figures in the 0.07--0.10
range. The heterogeneity puzzle (\emph{I}² = 62.4\%) remains largely
unresolved, with within-study variance (Level 2, 54.1\%) dominating
between-study variance (Level 3, 8.4\%), suggesting that DOI
operationalization and performance measure choices within papers are
more important sources of heterogeneity than cross-country institutional
differences.

\textbf{Limitations.} Five inferential constraints bound these
conclusions. First, the three moderator hypotheses are not confirmed in
the present \emph{k} = 238 corpus; this may reflect insufficient
regime-level \emph{k} (Frontier: \emph{k} = 3; SIDS: \emph{k} = 0)
rather than genuine null effects. A formal WoS/Scopus search targeting
Frontier and SIDS contexts is required before the ICRV null result can
be interpreted as definitive. Second, the Frontier-group anomaly
(\emph{k} = 3, \emph{r̄} = 0.349, dominated by Pouresmaeili et al.~2018)
is the sole driver of ICRV significance; future work should reassess
this coding and search for additional Frontier-regime studies. Third,
cDAI is measured at the study level using country-year World Bank
Digital Adoption Index scores (primary) or ITU Digital Development Index
(secondary), not direct replication of the WBES firm-level indicators
used in the companion primary studies; any systematic error in the
country-year DAI assignment, for instance, from vintage mismatch between
the study's data period and the DAI score's reference year, may
attenuate the moderation coefficient. Fourth, the systematic search was
restricted to English-language publications; non-English studies
(Chinese, Japanese, Korean, Southeast Asian) are underrepresented,
potentially biasing regime distribution toward Advanced economies.
Fifth, reliability is established on a 20\% double-coded subsample
(\emph{k} = 47 studies), with the two authors coding independently and
agreement reported in Table 3.1; the remaining 80\% is single-coded by
the first author against the codebook, which, while standard for
resource-bounded syntheses, cannot be independently validated.

\textbf{Future research.} Three directions follow directly from the null
moderator findings. First, a formal WoS/Scopus systematic search
targeting Frontier and SIDS economy I to P studies (targeted search
strings, Appendix B) would expand Frontier \emph{k} from 3 to
potentially 20--40 studies, enabling a meaningful test of the ICRV
Frontier cell. Second, a pre-registered replication of the ICRV, cDAI,
and DPL moderator hypotheses, with explicit power analysis for each
regime cell and OSF registration, would provide a definitive test of
whether these moderators reach significance in a \emph{k} ≥ 400 corpus.
Third, the substantial publication bias finding (\emph{k} = 57 imputed,
adj. \emph{r} = 0.035) motivates a dedicated publication bias audit:
surveying I to P researchers about unpublished null results, and
meta-analyzing grey literature and dissertations, would bound the true
population effect more precisely than trim-and-fill alone.

\hypertarget{funding}{%
\subsection{Funding}\label{funding}}

This research received no specific grant from any funding agency in the
public, commercial, or not-for-profit sectors.

\hypertarget{competing-interests}{%
\subsection{Competing Interests}\label{competing-interests}}

The authors declare no conflict of interest.

\hypertarget{data-availability}{%
\subsection{Data Availability}\label{data-availability}}

The coded study database (\texttt{forest\_data.csv}), R analysis scripts
(\texttt{metafor}), and the OSF registration protocol are available from
the corresponding author upon reasonable request. The PRISMA 2020
checklist is provided as supplementary material.

\hypertarget{declaration-of-generative-ai-and-ai-assisted-technologies-in-the-writing-process}{%
\subsection{Declaration of Generative AI and AI-Assisted Technologies in
the Writing
Process}\label{declaration-of-generative-ai-and-ai-assisted-technologies-in-the-writing-process}}

During the preparation of this work, the authors used two software aids,
each under full human control and verification. First, a purpose-built
large-language-model-assisted tool (M-AIDA, using Anthropic Claude) was
used to assist effect-size extraction and statistical conversion from
the primary studies; every value proposed by the tool was independently
checked, corrected where necessary, and permanently locked by the
Principal Investigator before entry into the analysis database. Second,
Grammarly was used to correct spelling, grammar, and punctuation in the
authors' own text. No generative AI tool selected studies, made
eligibility decisions, performed the statistical analysis, or generated
any interpretive or written content; the authors reviewed all outputs
and take full responsibility for the content of the publication.

\hypertarget{references}{%
\subsection{References}\label{references}}

Aguinis, H., Dalton, D. R., Bosco, F. A., Pierce, C. A., \& Dalton, C.
M. (2011). Meta-analytic choices and judgment calls: Implications for
theory and research. \emph{Journal of Management, 37}(1), 5--38.

Arte, P., \& Larimo, J. (2022). Moderating influence of product
diversification on the internationalization-performance relationship:
Insights from meta-analysis. \emph{Journal of Business Research, 139},
1408--1423.

Barney, J. (1991). Firm resources and sustained competitive advantage.
\emph{Journal of Management, 17}(1), 99--120.

Bausch, A., \& Krist, M. (2007). The effect of context-related
moderators on the internationalization--performance relationship:
Evidence from meta-analysis. \emph{Management International Review,
47}(3), 319--347.

Begg, C. B., \& Mazumdar, M. (1994). Operating characteristics of a rank
correlation test for publication bias. \emph{Biometrics, 50}(4),
1088--1101.

Bharadwaj, A., El Sawy, O. A., Pavlou, P. A., \& Venkatraman, N. (2013).
Digital business strategy: Toward a next generation of insights.
\emph{MIS Quarterly, 37}(2), 471--482.

Borenstein, M., Hedges, L. V., Higgins, J. P. T., \& Rothstein, H. R.
(2021). \emph{Introduction to meta-analysis} (2nd ed.). Wiley.

Brynjolfsson, E., Rock, D., \& Syverson, C. (2021). The productivity
J-curve: How intangibles complement general purpose technologies.
\emph{American Economic Journal: Macroeconomics, 13}(1), 333--372.

Bustamante, C. V., Mingo, S., \& Matusik, S. F. (2022). Institutions,
digital capabilities and the internationalization of SMEs. \emph{Journal
of International Business Studies, 53}(3), 524--546.

Cheung, M. W.-L. (2014). Modeling dependent effect sizes with
three-level meta-analyses. \emph{Psychological Methods, 19}(2),
211--226.

Cohen, J. (1988). \emph{Statistical power analysis for the behavioral
sciences} (2nd ed.). Lawrence Erlbaum.

Cooper, H. (2010). \emph{Research synthesis and meta-analysis: A
step-by-step approach} (4th ed.). Sage.

David, P. A. (1990). The dynamo and the computer: An historical
perspective on the modern productivity paradox. \emph{American Economic
Review, 80}(2), 355--361.

Authors (2024). Details omitted for double-blind review.

Authors (2025). Details omitted for double-blind review.

Duval, S., \& Tweedie, R. (2000). Trim and fill: A simple
funnel-plot-based method of testing and adjusting for publication bias
in meta-analysis. \emph{Biometrics, 56}(2), 455--463.

Egger, M., Smith, G. D., Schneider, M., \& Minder, C. (1997). Bias in
meta-analysis detected by a simple, graphical test. \emph{BMJ,
315}(7109), 629--634.

Hedges, L. V., \& Olkin, I. (1985). \emph{Statistical methods for
meta-analysis}. Academic Press.

Helpman, E., Melitz, M. J., \& Yeaple, S. R. (2004). Export versus FDI
with heterogeneous firms. \emph{American Economic Review, 94}(1),
300--316.

Jacquemin, A. P., \& Berry, C. H. (1979). Entropy measure of
diversification and corporate growth. \emph{Journal of Industrial
Economics, 27}(4), 359--369.

Johanson, J., \& Vahlne, J.-E. (1977). The internationalization process
of the firm: A model of knowledge development and increasing foreign
market commitments. \emph{Journal of International Business Studies,
8}(1), 23--32.

Katz, R., \& Callorda, F. (2018). \emph{The economic contribution of
broadband, digitization and ICT regulation}. ITU Telecommunication
Development Sector.

Khanna, T., \& Palepu, K. G. (2010). \emph{Winning in emerging markets:
A road map for strategy and execution}. Harvard Business Press.

Kirca, A. H., Hult, G. T. M., Deligonul, S., Perryy, M. Z., \& Cavusgil,
S. T. (2012). A multilevel examination of the drivers of firm
multinationality: A meta-analysis. \emph{Journal of Management, 38}(2),
502--530.

Kraus, S., Breier, M., Lim, W. M., Dabić, M., Kumar, S., Kanbach, D.,
Mukherjee, D., Corvello, V., Piñeiro-Chousa, J., Liguori, E.,
Palacios-Marqués, D., Schiavone, F., Ferraris, A., Fernandes, C., \&
Ferreira, J. J. (2022). Literature reviews as independent studies:
Guidelines for academic practice. \emph{Review of Managerial Science,
16}(8), 2577--2595.

Kogut, B., \& Zander, U. (1993). Knowledge of the firm and the
evolutionary theory of the multinational corporation. \emph{Journal of
International Business Studies, 24}(4), 625--645.

Landis, J. R., \& Koch, G. G. (1977). The measurement of observer
agreement for categorical data. \emph{Biometrics, 33}(1), 159--174.

Marano, V., Arregle, J.-L., Hitt, M. A., Spadafora, E., \& van Essen, M.
(2016). Home country institutions and the
internationalization-performance relationship: A meta-analytic review.
\emph{Journal of Management, 42}(5), 1075--1110.

North, D. C. (1990). \emph{Institutions, institutional change and
economic performance}. Cambridge University Press.

Orwin, R. G. (1983). A fail-safe N for effect size in meta-analysis.
\emph{Journal of Educational Statistics, 8}(2), 157--159.

Page, M. J., McKenzie, J. E., Bossuyt, P. M., Boutron, I., Hoffmann, T.
C., Mulrow, C. D., Shamseer, L., Tetzlaff, J. M., Akl, E. A., Brennan,
S. E., Chou, R., Glanville, J., Grimshaw, J. M., Hróbjartsson, A., Lalu,
M. M., Li, T., Loder, E. W., Mayo-Wilson, E., McDonald, S., \ldots{}
Moher, D. (2021). The PRISMA 2020 statement: An updated guideline for
reporting systematic reviews. \emph{BMJ, 372}, n71.
\url{https://doi.org/10.1136/bmj.n71}

Peng, M. W., Wang, D. Y. L., \& Jiang, Y. (2008). An institution-based
view of international business strategy: A focus on emerging economies.
\emph{Journal of International Business Studies, 39}(5), 920--936.

Peterson, R. A., \& Brown, S. P. (2005). On the use of beta coefficients
in meta-analysis. \emph{Journal of Applied Psychology, 90}(1), 175--181.

Raudenbush, S. W., \& Bryk, A. S. (2002). \emph{Hierarchical linear
models: Applications and data analysis methods} (2nd ed.). Sage.

Rosenthal, R. (1991). \emph{Meta-analytic procedures for social
research} (rev. ed.). Sage.

Rosenthal, R. (1994). Parametric measures of effect size. In H. Cooper
\& L. V. Hedges (Eds.), \emph{The handbook of research synthesis}
(pp.~231--244). Sage.

Schwens, C., Zapkau, F. B., Brouthers, K. D., \& Hollender, L. (2018).
Limits to outsourcing: A meta-analysis and empirical investigation.
\emph{Journal of International Business Studies, 49}(6), 682--703.

Stanley, T. D., \& Doucouliagos, H. (2014). Meta-regression
approximations to reduce publication selection bias. \emph{Research
Synthesis Methods, 5}(1), 60--78.

Stallkamp, M., \& Schotter, A. P. J. (2021). Platforms without borders?
The international strategies of digital platform firms. \emph{Global
Strategy Journal, 11}(1), 58--80.

Sutton, A. J., \& Higgins, J. P. T. (2008). Recent developments in
meta-analysis. \emph{Statistics in Medicine, 27}(5), 625--650.

Van den Noortgate, W., López-López, J. A., Marín-Martínez, F., \&
Sánchez-Meca, J. (2013). Three-level meta-analysis of dependent effect
sizes. \emph{Behavior Research Methods, 45}(2), 576--594.

Verhoef, P. C., Broekhuizen, T., Bart, Y., Bhattacharya, A., Qi Dong,
J., Fabian, N., \& Haenlein, M. (2021). Digital transformation: A
multidisciplinary reflection and research agenda. \emph{Journal of
Business Research, 122}, 889--901.

Vevea, J. L., \& Woods, C. M. (2005). Publication bias in research
synthesis: Sensitivity analysis using a priori weight functions.
\emph{Psychological Methods, 10}(4), 428--443.

Viechtbauer, W. (2010). Conducting meta-analyses in R with the metafor
package. \emph{Journal of Statistical Software, 36}(3), 1--48.

Wernerfelt, B. (1984). A resource-based view of the firm.
\emph{Strategic Management Journal, 5}(2), 171--180.

World Bank. (2023). \emph{Worldwide Governance Indicators}.
\url{https://info.worldbank.org/governance/wgi/}

Wu, J., Wang, C., Hong, J., Piperopoulos, P., \& Zhuo, S. (2022).
Internationalization and innovation performance of emerging market
enterprises: The role of host-country institutional development.
\emph{Journal of World Business, 52}(2), 192--203.

Zaheer, S. (1995). Overcoming the liability of foreignness.
\emph{Academy of Management Journal, 38}(2), 341--363.

\emph{Appendices A--C (PRISMA flow, coding protocol,
MetaEssentials/\texttt{metafor} consistency check) are provided as
online supplementary material.}

\end{document}
